\chapter{Étude Préliminaire et Analyse des Besoins}

\section*{Introduction}
\addcontentsline{toc}{section}{Introduction}

Dans ce premier chapitre, nous présentons l'Institut National Supérieur des Sciences et Techniques d'Abéché (INSTA), notre établissement de formation. Nous exposons ensuite le contexte général du projet, l'analyse de l'existant, la problématique identifiée, ainsi que les besoins fonctionnels et non fonctionnels du système à concevoir.

%=============================================
\section{Contexte Institutionnel : l'INSTA}
%=============================================

\subsection{Présentation générale}

L'Institut National Supérieur des Sciences et Techniques d'Abéché (INSTA) est un établissement public d'enseignement supérieur à caractère scientifique et technique, situé dans la ville d'Abéché (province du Ouaddaï), au Nord-Est du Tchad. Créé initialement sous le nom d'Institut Universitaire des Sciences et Techniques d'Abéché (IUSTA) en 1997, il a été rebaptisé INSTA par l'Ordonnance N°004/PR/MES/2015 du 02 mars 2015.

\subsection{Missions de l'INSTA}

L'institut poursuit plusieurs missions fondamentales :

\begin{itemize}
    \item \textbf{Formation académique et professionnelle} : former des étudiants dans divers domaines techniques et scientifiques, en vue de répondre aux besoins du marché de l'emploi.
    \item \textbf{Recherche et innovation} : promouvoir la recherche scientifique et valoriser les résultats pour répondre aux besoins du pays.
    \item \textbf{Partenariats} : établir des collaborations avec des institutions nationales et internationales pour renforcer la qualité de l'enseignement.
    \item \textbf{Contribution au développement} : fournir au marché de l'emploi des ingénieurs compétents capables d'apporter des solutions aux problématiques locales.
\end{itemize}

\subsection{Offres de formation}

L'INSTA offre deux cycles principaux :

\begin{itemize}
    \item \textbf{Premier cycle} : Licence professionnelle (3 ans) formant des licenciés professionnels dans les domaines scientifiques et techniques.
    \item \textbf{Deuxième cycle} : Master (2 ans), divisé en Master professionnel et recherche, dans les domaines de Génie Électrique, Génie Mécanique, Génie Énergétique, Sciences Biomédicales et Sciences et Techniques d'Élevage.
\end{itemize}

\subsection{Départements de l'INSTA}

L'institut est composé de plusieurs départements :

\begin{itemize}
    \item Département de Génie Informatique (GI) — option Génie Logiciel
    \item Département de Génie Mécanique (GM)
    \item Département de Génie Électrique (GE)
    \item Département de Génie Énergétique (GEN)
    \item Département de Télécommunication et Réseau (TLC)
    \item Département Conception, Installation et Maintenance des Systèmes (CIMS)
    \item Département des Sciences et Techniques d'Élevage (STE)
    \item Département des Sciences Biomédicales et Pharmaceutiques (SBMP)
\end{itemize}

%=============================================
\section{Contexte Général du Projet}
%=============================================

\subsection{Le Génie Logiciel et la modélisation UML}

Le Génie Logiciel est la discipline qui s'intéresse à la conception, au développement, à la maintenance et à la gestion de systèmes logiciels de qualité. Parmi les pratiques fondamentales de cette discipline, la modélisation occupe une place centrale. Elle permet de représenter un système sous forme abstraite avant son implémentation, facilitant ainsi la communication entre les parties prenantes et la détection précoce des erreurs de conception.

L'UML (\textit{Unified Modeling Language}) est le langage de modélisation le plus utilisé dans l'industrie du logiciel. Il propose un ensemble de diagrammes standardisés permettant de décrire l'architecture statique et dynamique d'un système logiciel. Pour notre projet, nous avons utilisé trois types de diagrammes :

\begin{table}[htbp]
\centering
\renewcommand{\arraystretch}{1.4}
\begin{tabular}{|l|p{5cm}|l|}
\hline
\rowcolor{gray!30}
\textbf{Diagramme} & \textbf{Objectif} & \textbf{Type} \\
\hline
Diagramme de classes & Représenter la structure statique du système : classes, attributs, méthodes et relations & Statique \\
\hline
Diagramme de cas d'utilisation & Décrire les interactions entre les acteurs et le système & Fonctionnel \\
\hline
Diagramme de séquence & Illustrer les échanges entre composants dans un ordre temporel & Dynamique \\
\hline
\end{tabular}
\caption{Diagrammes UML utilisés dans le projet}
\end{table}

\subsection{Le Développement Dirigé par les Modèles (MDD)}

Le \textbf{Développement Dirigé par les Modèles} (MDD — \textit{Model-Driven Development}) est un paradigme de développement logiciel qui place le modèle au cœur du processus de développement. L'idée fondatrice du MDD est que le modèle n'est plus un simple document de documentation, mais un artefact de première classe à partir duquel le code source peut être généré automatiquement.

Dans ce paradigme, on distingue :
\begin{itemize}
    \item Le \textbf{PIM} (\textit{Platform Independent Model}) : le modèle UML indépendant de toute plateforme
    \item Le \textbf{PSM} (\textit{Platform Specific Model}) : le code source généré pour une plateforme cible (Python, PHP, etc.)
    \item La \textbf{transformation} : le processus automatique qui convertit le PIM en PSM
\end{itemize}

Notre projet s'inscrit exactement dans cette approche : le diagramme UML est le PIM, et le code Python ou PHP généré est le PSM.

\subsection{L'Intelligence Artificielle générative}

L'avènement des \textbf{Grands Modèles de Langage} (LLM — \textit{Large Language Models}) représente une révolution dans le domaine de l'Intelligence Artificielle. Ces modèles, entraînés sur des milliards de lignes de texte et de code source, sont capables de comprendre et de générer du code source de manière contextuelle et cohérente.

Dans le cadre de ce projet, nous exploitons le modèle \textbf{Groq Llama 3.3-70b}, un LLM open source de 70 milliards de paramètres développé par Meta et hébergé sur l'infrastructure d'inférence ultra-rapide \textbf{Groq}, qui permet des temps de génération de l'ordre de 2 à 5 secondes.

\subsection{Le format XMI}

Le format \textbf{XMI} (\textit{XML Metadata Interchange}) est un standard défini par l'OMG\footnote{Object Management Group : \reporturl{https://www.omg.org}} pour l'échange de modèles UML entre différents outils de modélisation. C'est un fichier texte structuré en XML qui contient la description complète d'un modèle UML : classes, attributs, méthodes, relations et héritage.

StarUML\fnStarUML, l'outil de modélisation que nous utilisons, permet d'exporter les diagrammes de classes au format XMI 2.1. Notre application parse ce fichier pour en extraire automatiquement la structure orientée objet et générer le code correspondant.

%=============================================
\section{Analyse de l'Existant}
%=============================================

\subsection{Outils existants de transformation UML vers code}

Plusieurs outils de transformation UML vers code existent sur le marché. Nous en avons analysé les principaux afin d'identifier leurs limites :

\begin{table}[htbp]
\centering
\renewcommand{\arraystretch}{1.4}
\begin{tabular}{|l|l|l|p{4.5cm}|}
\hline
\rowcolor{gray!30}
\textbf{Outil} & \textbf{Langages générés} & \textbf{Format entrée} & \textbf{Limites principales} \\
\hline
Acceleo (Eclipse) & Java, Python, PHP & XMI/EMF & Très complexe, lourd à installer \\
\hline
AndroMDA & Java EE uniquement & XMI & Limité à Java, pas d'IA \\
\hline
GenMyModel & Java, Python & En ligne & Payant, pas open source \\
\hline
StarUML Plugins & JavaScript, Java & .mdj & Limité, pas d'IA générative \\
\hline
Solutions manuelles & Tous langages & Tous formats & Lent, source d'erreurs \\
\hline
\end{tabular}
\caption{Comparatif des outils existants de transformation UML vers code}
\end{table}

\subsection{Limites des solutions existantes}

L'analyse comparative révèle plusieurs lacunes communes :

\begin{itemize}
    \item \textbf{Complexité d'installation} : des outils comme Acceleo nécessitent une installation lourde et une configuration complexe
    \item \textbf{Absence d'IA} : aucun outil n'intègre de modèle d'IA générative pour produire un code contextuel
    \item \textbf{Formats d'entrée limités} : la plupart n'acceptent que le format XMI, excluant les images et les PDF
    \item \textbf{Coût} : certaines solutions performantes sont payantes et propriétaires
    \item \textbf{Pas de génération de serveur} : aucun outil ne génère automatiquement un serveur REST opérationnel
\end{itemize}

\subsection{Valeur ajoutée de notre solution}

Face à ces limitations, \textbf{UML-to-Code} apporte une réponse innovante :
\begin{itemize}
    \item Intégration d'une IA générative (\textbf{Groq Llama 3.3-70b}) pour un code contextuel et complet
    \item Support de plusieurs formats d'entrée : XMI, image PNG/JPG, PDF
    \item Génération automatique d'un serveur Flask complet avec routes CRUD
    \item Interface web moderne avec Monaco Editor, historique local, mode sombre/clair
    \item Export en archive ZIP du projet complet prêt à déployer
    \item Déploiement en ligne accessible via Vercel\fnVercel{} et Render\fnRender
    \item Solution légère, open source et sans installation complexe
\end{itemize}

%=============================================
\section{Spécification des Besoins}
%=============================================

\subsection{Acteurs du système}

\begin{table}[htbp]
\centering
\renewcommand{\arraystretch}{1.4}
\begin{tabular}{|l|p{4cm}|p{6cm}|}
\hline
\rowcolor{gray!30}
\textbf{Acteur} & \textbf{Rôle} & \textbf{Interactions principales} \\
\hline
\textbf{Utilisateur} & Étudiant ou développeur utilisant l'outil & Upload fichier, choix options, génération, visualisation Monaco Editor, copie, téléchargement, consultation historique \\
\hline
\textbf{API Groq (IA)} & Modèle Llama 3.3-70b & Reçoit la description UML extraite, retourne le code généré \\
\hline
\textbf{Backend Flask} & Serveur Python & Orchestration, parsing XMI/image/PDF, appel API Groq, génération ZIP \\
\hline
\textbf{localStorage} & Stockage navigateur & Persiste l'historique des 10 dernières générations et le thème \\
\hline
\end{tabular}
\caption{Acteurs du système UML-to-Code}
\end{table}

\subsection{Exigences Fonctionnelles}

Les exigences fonctionnelles sont regroupées par module dans le tableau~\ref{tab:ef-complet}. Les priorités \og Haute \fg{} et \og Moyenne \fg{} guident l'ordre de réalisation des incréments.

\begin{longtable}{|c|l|P{6.8cm}|c|}
\hline
\rowcolor{blue!12}
\textbf{ID} & \textbf{Module} & \textbf{Exigence} & \textbf{Priorité} \\
\hline
\endfirsthead
\hline
\rowcolor{blue!12}
\textbf{ID} & \textbf{Module} & \textbf{Exigence} & \textbf{Priorité} \\
\hline
\endhead
EF-01 & Upload & Accepter les fichiers XMI exportés depuis StarUML\fnStarUML & Haute \\
\hline
EF-02 & Upload & Accepter les images PNG et JPG de diagrammes UML & Haute \\
\hline
EF-03 & Upload & Accepter les fichiers PDF contenant un diagramme UML & Moyenne \\
\hline
EF-04 & Upload & Valider le format du fichier avant traitement & Haute \\
\hline
EF-05 & Upload & Afficher un message d'erreur en cas de format invalide & Haute \\
\hline
EF-06 & Upload & Afficher un aperçu visuel du diagramme pour PNG/JPG & Moyenne \\
\hline
EF-07 & Parsing & Parser les classes, attributs et méthodes depuis un XMI & Haute \\
\hline
EF-08 & Parsing & Extraire les relations d'héritage entre classes & Haute \\
\hline
EF-09 & Parsing & Extraire les associations et agrégations & Moyenne \\
\hline
EF-10 & Parsing & Envoyer l'image/PDF à Groq (vision IA)\fnGroq & Haute \\
\hline
EF-11 & Génération & Appeler l'API Groq (Llama 3.3-70b) avec la description UML & Haute \\
\hline
EF-12 & Génération & Générer le code Python (\texttt{\_\_init\_\_}, getters, setters) & Haute \\
\hline
EF-13 & Génération & Générer le code PHP (constructeur, getters) & Haute \\
\hline
EF-14 & Génération & Générer un serveur Flask avec routes CRUD\fnFlask & Haute \\
\hline
EF-15 & Génération & Choix du langage : Python ou PHP & Haute \\
\hline
EF-16 & Génération & Choix du mode : Algorithmique ou IA Groq & Haute \\
\hline
EF-17 & Interface & Afficher le code dans Monaco Editor\fnMonaco & Haute \\
\hline
EF-18 & Interface & Onglets Code et Serveur Flask séparés & Haute \\
\hline
EF-19 & Interface & Copie du code en un clic & Haute \\
\hline
EF-20 & Interface & Téléchargement du fichier \texttt{.py} ou \texttt{.php} & Haute \\
\hline
EF-21 & Interface & Export en archive ZIP & Haute \\
\hline
EF-22 & Interface & Barre de progression en temps réel & Haute \\
\hline
EF-23 & Interface & Historique local (10 générations) & Haute \\
\hline
EF-24 & Interface & Mode sombre / mode clair & Moyenne \\
\hline
EF-25 & Interface & Notifications toast & Moyenne \\
\hline
\caption{Exigences fonctionnelles complètes du système UML-to-Code}
\label{tab:ef-complet}
\end{longtable}

\subsection{Exigences Non Fonctionnelles}

\begin{table}[htbp]
\centering
\renewcommand{\arraystretch}{1.15}
\begin{tabular}{|l|P{4.8cm}|P{4.2cm}|}
\hline
\rowcolor{blue!12}
\textbf{Catégorie} & \textbf{Exigence} & \textbf{Critère mesurable} \\
\hline
Performance & Temps de réponse de la génération & $<$ 15 secondes par diagramme \\
\hline
Disponibilité & Application accessible en ligne\fnVercel & 24h/24 via Vercel \\
\hline
Fiabilité & Taux de succès de la génération & $>$ 90\% sur les cas de test \\
\hline
Maintenabilité & Code source documenté et structuré & Commentaires sur chaque fonction \\
\hline
Portabilité & Fonctionne sur Windows et Linux & Testé sur les deux systèmes \\
\hline
Sécurité & Clé API en variable d'environnement & Fichier \texttt{.env} non versionné \\
\hline
Utilisabilité & Interface intuitive sans formation & Prise en main $<$ 5 minutes \\
\hline
Responsivité & Interface adaptée mobile et desktop & 320px à 1440px \\
\hline
\end{tabular}
\caption{Exigences non fonctionnelles}
\label{tab:enf}
\end{table}

\section*{Conclusion}
\addcontentsline{toc}{section}{Conclusion}

Dans ce premier chapitre, nous avons présenté l'établissement d'accueil INSTA et exposé le contexte général du projet. Nous avons analysé les solutions existantes de transformation UML vers code et identifié leurs limites, ce qui nous a permis de justifier la valeur ajoutée de notre solution UML-to-Code\fnGitHub. Enfin, nous avons spécifié les besoins fonctionnels et non fonctionnels du système à concevoir. Le chapitre suivant sera consacré à la conception et la modélisation du système.